To enforce this discrete precedence order at the continuous control level, we embed it into a responsibility-aware control barrier function (CBF)~\cite{cosner2023learning}. 
We define a directional indicator $\psi_{ij} \in \{-1, 1\}$ based on the yielding order and allocate an asymmetric avoidance responsibility $\gamma_{ij} = \rho \psi_{ij}$.
The yielding agent takes on a tighter constraint within the local quadratic program (QP):
\begin{equation}
    L_{g_i}h_{ij}u_i + \frac{1}{2}(L_{f_{ij}}h_{ij} + \alpha(h_{ij})) - \gamma_{ij} \ge 0.
\end{equation} 
By exchanging only the scalar $\mu_i$, agents dynamically reallocate avoidance efforts without sharing future trajectories, strictly preserving pairwise collision avoidance. 

% The energy quantities in Sec.~\ref{sec:4.B} determine the discrete precedence, while the continuous control layer enforces the associated pairwise collision-avoidance constraints.
% For an interacting pair $(i,j)$, define the collision-avoidance barrier
% \begin{equation}
%     h_{ij}(x_i,x_j):=\|p_i-p_j\|^2-d_{\min}^2,
%     \label{eq:collision_barrier}
% \end{equation}
% where $d_{\min}>0$ is the required separation distance.
% Assuming relative degree one, the shared CBF condition is
% \begin{equation}
%     L_{f_{ij}}h_{ij}
%     +
%     L_{g_i}h_{ij}u_i
%     +
%     L_{g_j}h_{ij}u_j
%     +
%     \alpha_c(h_{ij})
%     \geq0.
%     \label{eq:shared_collision_cbf}
% \end{equation}
% We encode the precedence in Eq.~\eqref{eq:energy_precedence} through the directional variable
% \begin{equation}
%     \psi_{ij}
%     :=
%     \begin{cases}
%         1, & i\rightarrow j,\\
%         -1, & j\rightarrow i,
%     \end{cases}
%     \qquad
%     \psi_{ji}=-\psi_{ij}.
%     \label{eq:precedence_direction}
% \end{equation}
% Let $\rho_{ij}=\rho_{ji}>0$ denote the responsibility magnitude and set
% \begin{equation}
%     \gamma_{ij}:=\rho_{ij}\psi_{ij},
%     \qquad
%     \gamma_{ji}:=-\gamma_{ij}.
%     \label{eq:complementary_responsibility_terms}
% \end{equation}
% Because the local constraint contains $-\gamma_{ij}$, the yielding robot receives the tighter constraint and hence greater collision-avoidance responsibility.

% Let $\mathcal{N}_i^{\mathrm{act}}$ denote the robots currently interacting with robot $i$.
% Robot $i$ computes its control input from
% \begin{equation}
% \begin{aligned}
%     u_i^\star:=
%     \arg&\min_{u_i\in\mathcal U_i}\quad
%     \frac{1}{2}
%     \|u_i-u_i^{\mathrm{nom}}\|_{W_i}^{2}\\
%     \mathrm{s.t.}&\quad
%     L_{g_i}h_{ij}u_i
%     +
%     \frac{1}{2}
%     \left(
%         L_{f_{ij}}h_{ij}
%         +
%         \alpha_c(h_{ij})
%     \right)
%     -
%     \gamma_{ij}
%     \geq0,\\
%     &\hspace{52.5mm}
%     \forall j\in\mathcal N_i^{\mathrm{act}},
% \end{aligned}
% \label{eq:local_precedence_qp}
% \end{equation}
% where $u_i^{\mathrm{nom}}$ is the task-level reference input and $W_i\succ0$.
% For each pair $(i,j)$, summing the two local constraints and using $\gamma_{ij}+\gamma_{ji}=0$ recovers the standard pairwise CBF condition.
% Thus, the allocation changes which robot contributes more to collision avoidance without changing the underlying safety requirement.
% Each robot simultaneously enforces all active pairwise constraints in Eq.~\eqref{eq:local_precedence_qp}.
% Under standard CBF regularity and QP-feasibility conditions, the collision-free sets are forward invariant.
% The resulting controller therefore satisfies the safety-preserving precedence-integration requirement of Problem~\ref{prob:main}.